% Math/Theory Summary Template (Mode B)
% For pure math explorations without data pipelines.
% Goal: self-contained document that a reader with fresh eyes can follow.
%
% Key principles:
%   - State the question and answer upfront (Introduction)
%   - List all proof steps before diving into details (Proof overview)
%   - Each step: Lemma statement -> Proof -> Implication (what this buys us)
%   - Discussion: what the argument does NOT use, remaining gaps
%   - Use model primitives directly in equations; avoid auxiliary notation

\documentclass[12pt]{article}

% --- Packages ---
\usepackage{lmodern}
\usepackage[T1]{fontenc}
\usepackage{amsmath,amssymb,amsfonts,mathtools,amsthm}
\usepackage{booktabs}
\usepackage{hyperref}
\usepackage{xcolor}
\usepackage{enumitem}
\usepackage{setspace}

% --- Page Layout ---
\usepackage[margin=1in]{geometry}
\setlength{\parindent}{0pt}
\setlength{\parskip}{0.8em}

% --- Section Formatting ---
\usepackage{titlesec}
\titleformat{\section}{\normalfont\scshape\centering}{\thesection.}{0.5em}{}
\titleformat{\subsection}{\normalfont\bfseries}{\thesubsection.}{0.5em}{}

% --- Hyperlinks ---
\hypersetup{
  colorlinks   = true,
  urlcolor     = blue,
  linkcolor    = blue,
  citecolor    = black,
}

% --- Theorem environments ---
\newtheorem{theorem}{Theorem}
\newtheorem{lemma}[theorem]{Lemma}
\newtheorem{corollary}[theorem]{Corollary}
\newtheorem{definition}[theorem]{Definition}

% --- Common Math Shortcuts (add project-specific ones as needed) ---
\newcommand{\E}{\mathbb{E}}
\DeclareMathOperator{\supp}{supp}

\begin{document}

% --- Title Block ---
\begin{center}
{\Large \textbf{PLACEHOLDER: Title}}

{\large PLACEHOLDER: Subtitle}

\medskip

Internal Report \hspace{1em}|\hspace{1em} \today
\end{center}


%-----------------------------------------------------%
\section{Introduction}\label{sec:intro}
%-----------------------------------------------------%

% State the question precisely. What is the reader about to learn?
% PLACEHOLDER: "This note asks: [precise question]?"

% State the answer. Don't make the reader wait.
% PLACEHOLDER: "The answer is [yes/no]. [One-paragraph summary of the result.]"

% Sketch the proof strategy in 2-3 sentences.
% PLACEHOLDER: "The proof proceeds by [high-level strategy]."

% Flag any remaining gaps honestly.
% PLACEHOLDER: \textcolor{red}{One gap remains: [description].}


%-----------------------------------------------------%
\section{Main Result}\label{sec:result}
%-----------------------------------------------------%

% State the theorem formally. This is the reference point for everything below.

% \begin{theorem}[PLACEHOLDER: descriptive name]\label{thm:main}
% PLACEHOLDER: formal statement.
% \end{theorem}


%-----------------------------------------------------%
\section{Proof}\label{sec:proof}
%-----------------------------------------------------%

% List ALL proof steps upfront as a numbered list.
% This gives the reader the full architecture before diving into details.
% Each step should be one sentence.

% The proof has N steps:
% \begin{enumerate}
%   \item PLACEHOLDER: Step 1 description.
%   \item PLACEHOLDER: Step 2 description.
%   \item ...
% \end{enumerate}


% --- Step 1 ---
% \subsection{Step 1: PLACEHOLDER title.}\label{sec:step1}
%
% \begin{lemma}\label{lem:step1}
% PLACEHOLDER: formal statement.
% \end{lemma}
%
% \begin{proof}
% PLACEHOLDER: proof body.
% \end{proof}
%
% \textit{Implication:} PLACEHOLDER: what this step buys us for the overall argument.
% Why does this matter? What would break without it?


% --- Step 2 ---
% \subsection{Step 2: PLACEHOLDER title.}\label{sec:step2}
% [Same pattern: Lemma -> Proof -> Implication]


% --- Combining the steps ---
% \subsection{Combining the steps.}\label{sec:combine}
%
% Restate the proof skeleton with cross-references:
% Take any [object]. By Lemma~\ref{lem:step1}, [consequence].
% By Lemma~\ref{lem:step2}, [consequence]. ...
% Therefore [main result]. \qed


%-----------------------------------------------------%
\section{Discussion}\label{sec:discussion}
%-----------------------------------------------------%

% What the argument does NOT use.
% This clarifies the generality of the result.
% PLACEHOLDER: "The proof never examines [aspect]. It works for any [object] because: ..."

% Remaining gaps.
% PLACEHOLDER: "[Description of what is assumed but not proved.]"


%-----------------------------------------------------%
% OPTIONAL: Appendix for alternative proofs or extensions
%-----------------------------------------------------%

% \appendix
% \section{PLACEHOLDER: Alternative Argument}\label{app:alt}
% Use when there is a second proof technique worth recording,
% or an extension that doesn't fit the main flow.


\end{document}
